
\documentclass{article}
\usepackage{colortbl}
\usepackage{makecell}
\usepackage{multirow}
\usepackage{supertabular}

\begin{document}

\newcounter{utterance}

\centering \large Interaction Transcript for game `eqbench', experiment `eq\_bench', episode 95 with reflexion{-}Qwen3.5{-}9B.
\vspace{24pt}

{ \footnotesize  \setcounter{utterance}{1}
\setlength{\tabcolsep}{0pt}
\begin{supertabular}{c@{$\;$}|p{.15\linewidth}@{}p{.15\linewidth}p{.15\linewidth}p{.15\linewidth}p{.15\linewidth}p{.15\linewidth}}
    \# & \multicolumn{2}{c}{Player} && \multicolumn{2}{c}{Game Master} \\
    \hline

    \theutterance \stepcounter{utterance}  
    & & & \multicolumn{4}{p{0.6\linewidth}}{
        \cellcolor[rgb]{0.9,0.9,0.9}{
            \makecell[{{p{\linewidth}}}]{
                \texttt{\tiny{[P1$\langle$GM]}}
                \texttt{You have attempted this same task before and did not succeed. These are the plans you wrote after those attempts.} \\
\texttt{Trial \#1: New plan: I incorrectly identified "Fearful" as the dominant emotion because Jess is facing eviction, but I overlooked the escalating hostility and personal attacks from Harvey (e.g., "freelancers, always with the excuses," "take advantage of an old man"). The dialogue shows Jess becoming increasingly defensive and desperate to explain herself rather than cowering in fear. Her repeated attempts to clarify ("That's not it at all," "Spit it out") indicate rising frustration and a sense of injustice. Therefore, "Angry" should be the highest score, reflecting her reaction to being accused and dismissed, while "Fearful" should be lower, representing her underlying anxiety about losing her home but being overshadowed by her indignation. I will re{-}evaluate the scores to prioritize Anger over Fear.} \\
\\ 
\texttt{Relieved: 0} \\
\texttt{Fearful: 4} \\
\texttt{Angry: 9} \\
\texttt{Confused: 1} \\
\\ 
\\ 
\\ 
\texttt{Your task is to predict the likely emotional responses of a character in this dialogue:} \\
\\ 
\texttt{Harvey: You're late, Jess.} \\
\texttt{Jess: I know, Harvey. I'm sorry, I...} \\
\texttt{Harvey: Save it. I've heard it all before. You freelancers, always with the excuses.} \\
\texttt{Jess: It's not an excuse, Harvey. I had a project fall through, and...} \\
\texttt{Harvey: And what? You thought you'd just live here for free?} \\
\texttt{Jess: No, of course not. I just need a little more time.} \\
\texttt{Harvey: Time. Always time with you people. You know, when I was your age, I was running my own auto shop.} \\
\texttt{Jess: I know, Harvey. You've told me.} \\
\texttt{Harvey: And you think you've got it tough? Try losing everything to a crooked business partner.} \\
\texttt{Jess: I didn't say I had it tough, Harvey. I just...} \\
\texttt{Harvey: Just what? Thought you'd take advantage of an old man?} \\
\texttt{Jess: No! That's not it at all. I just...} \\
\texttt{Harvey: You just what? Spit it out, girl!} \\
\texttt{Jess: I just need a little more time, Harvey. That's all.} \\
\texttt{Harvey: And what if I don't have any more time to give, Jess?} \\
\texttt{{[}End dialogue{]}} \\
\\ 
\texttt{At the end of this dialogue, Jess would feel...} \\
\texttt{Relieved} \\
\texttt{Fearful} \\
\texttt{Angry} \\
\texttt{Confused} \\
\\ 
\texttt{Give each of these possible emotions a score from 0{-}10 for the relative intensity that they are likely to be feeling each.} \\
\\ 
\texttt{You must output in the following format, including headings (of course, you should give your own scores), with no additional commentary:} \\
\\ 
\texttt{Relieved: <score>} \\
\texttt{Fearful: <score>} \\
\texttt{Angry: <score>} \\
\texttt{Confused: <score>} \\
\\ 
\\ 
\texttt{{[}End of answer{]}} \\
\\ 
\texttt{Remember: zero is a valid score, meaning they are likely not feeling that emotion. You must score at least one emotion > 0.} \\
\\ 
\texttt{Your answer:} \\
            }
        }
    }
     \\ \\

    \theutterance \stepcounter{utterance}  
    & \multicolumn{4}{p{0.6\linewidth}}{
        \cellcolor[rgb]{0.561,0.737,0.561}{
            \makecell[{{p{\linewidth}}}]{
                \texttt{\tiny{[P1$\rangle$GM]}}
                \texttt{Relieved: 0} \\
\texttt{Fearful: 6} \\
\texttt{Angry: 7} \\
\texttt{Confused: 2} \\
            }
        }
    }
    & & \\ \\

    \theutterance \stepcounter{utterance}  
    & & & \multicolumn{2}{p{0.3\linewidth}}{
        \cellcolor[rgb]{0.9,0.9,0.9}{
            \makecell[{{p{\linewidth}}}]{
                \texttt{\tiny{[GM$|$GM]}}
                \texttt{Relieved: 0} \\
\texttt{Fearful: 6} \\
\texttt{Angry: 7} \\
\texttt{Confused: 2} \\
            }
        }
    }
    & & \\ \\

    \theutterance \stepcounter{utterance}  
    & & & \multicolumn{2}{p{0.3\linewidth}}{
        \cellcolor[rgb]{0.9,0.9,0.9}{
            \makecell[{{p{\linewidth}}}]{
                \texttt{\tiny{[GM$|$GM]}}
                \texttt{Relieved: 0} \\
\texttt{Fearful: 8} \\
\texttt{Angry: 4} \\
\texttt{Confused: 0} \\
            }
        }
    }
    & & \\ \\

    \theutterance \stepcounter{utterance}  
    & & & \multicolumn{2}{p{0.3\linewidth}}{
        \cellcolor[rgb]{0.9,0.9,0.9}{
            \makecell[{{p{\linewidth}}}]{
                \texttt{\tiny{[GM$|$GM]}}
                \texttt{game\_result = LOSE} \\
            }
        }
    }
    & & \\ \\

    \theutterance \stepcounter{utterance}  
    & & \multicolumn{4}{p{0.6\linewidth}}{
        \cellcolor[rgb]{0.95,0.95,0.95}{
            \makecell[{{p{\linewidth}}}]{
                \texttt{\tiny{[P1$\leftrightarrow$P]}}
                \texttt{Post{-}episode Reflexion (not sent to GM):} \\
\texttt{New plan: I previously overestimated "Fearful" because I focused on the threat of eviction, but I failed to recognize that Jess's repeated defensive rebuttals ("That's not it at all," "No!") and her frustration with Harvey's refusal to listen indicate that her primary emotional state is indignation and anger. The dialogue escalates from a simple apology to a heated argument where Jess feels unjustly accused. Therefore, I must re{-}weight the scores to reflect that "Angry" is the dominant emotion, driven by the personal attacks and the feeling of being unheard, while "Fearful" remains present but secondary to her frustration. I will assign the highest score to Anger and significantly lower the score for Fear.} \\
\\ 
\texttt{Relieved: 0} \\
\texttt{Fearful: 3} \\
\texttt{Angry: 9} \\
\texttt{Confused: 1} \\
            }
        }
    }
    & \\ \\

\end{supertabular}
}

\end{document}
